\documentclass[aspectratio=169, table, xcdraw]{beamer}

% Tema y Colores Institucionales[cite: 3]
\usetheme{Madrid}
\usecolortheme{default}
\definecolor{ucbblue}{RGB}{0, 51, 102}

\setbeamercolor{structure}{fg=ucbblue}
\setbeamercolor*{palette primary}{fg=white,bg=ucbblue}
\setbeamercolor*{palette secondary}{fg=white,bg=ucbblue!80}
\setbeamercolor*{palette tertiary}{fg=white,bg=ucbblue!90}
\setbeamercolor{title}{fg=white, bg=ucbblue}
\setbeamercolor{frametitle}{fg=white, bg=ucbblue}
\setbeamercolor{item}{fg=ucbblue}

\usepackage[T1]{fontenc}
\usepackage[utf8]{inputenc}
\usepackage{tgpagella}
\usefonttheme{serif}
\usepackage[spanish, es-noshorthands]{babel}
\renewcommand{\tablename}{Tabla}
\renewcommand{\figurename}{Figura}

\usepackage{amsmath, amssymb, bm}
\usepackage{graphicx}
\usepackage{booktabs}
\usepackage[most]{tcolorbox}
\usepackage{tikz}
\usetikzlibrary{shapes.geometric, arrows.meta, positioning, calc}

\title[ABB IRB 120: Cinemática Directa]{Cinemática Directa mediante Parámetros D-H}
\subtitle{De la Convención D-H a la Pose del Efector Final}
\author[B. Quiroga Turdera]{M.Sc. Ing. Bernardo Quiroga Turdera}
\institute[UCB Tarija]{Universidad Católica Boliviana ``San Pablo'' — Sede Tarija \\ Robótica (IMT-342)}
\date{Martes, 1 de Septiembre de 2026}

\begin{document}

\begin{frame}
    \titlepage
\end{frame}

\begin{frame}{El Camino hacia el Hito 1}
    \begin{center}
    \begin{tikzpicture}[
        node distance=1cm,
        box/.style={rectangle, draw=ucbblue, thick, fill=white, text width=2.4cm, align=center, rounded corners, minimum height=1cm, font=\scriptsize},
        activebox/.style={rectangle, draw=ucbblue, thick, fill=ucbblue!15, text width=2.4cm, align=center, rounded corners, minimum height=1cm, font=\scriptsize\bfseries},
        arrow/.style={-Stealth, thick, ucbblue}
    ]
        \node[box] (f1) {1. Marcos D-H};
        \node[box, right=of f1] (f2) {2. Tabla D-H};
        \node[activebox, right=of f2] (f3) {3. Cinemática Directa Numérica};
        \node[box, below=of f3] (f4) {4. Modelo URDF};
        \node[box, left=of f4] (f5) {5. Validación RViz};
        \node[box, fill=green!10, left=of f5] (f6) {6. Defensa EC1};

        \draw[arrow] (f1) -- (f2);
        \draw[arrow] (f2) -- (f3);
        \draw[arrow] (f3) -- (f4);
        \draw[arrow] (f4) -- (f5);
        \draw[arrow] (f5) -- (f6);
    \end{tikzpicture}
    \end{center}
    \vspace{0.4cm}
    \begin{block}{Objetivo de la Sesión[cite: 1, 2]}
        Completar el paso pendiente: Transformar la tabla D-H en matrices individuales, componer la cadena cinemática y calcular matemáticamente la posición y orientación de la brida.
    \end{block}
\end{frame}

\begin{frame}{Cinemática Directa: Definición del Problema}
    La cinemática directa es una función que determina la pose del efector final a partir de las variables articulares y la geometría conocida del robot[cite: 1].
    
    \vspace{0.3cm}
    \begin{center}
        \large
        $\boxed{\mathbf{q} \longmapsto {}^{0}T_n(\mathbf{q})}$
    \end{center}
    \vspace{0.3cm}
    
    \textbf{Para un manipulador 6R (ej. ABB IRB 120):}
    \begin{itemize}
        \item \textbf{Entrada:} Configuración articular $\mathbf{q} = [q_1, q_2, q_3, q_4, q_5, q_6]^T$[cite: 1].
        \item \textbf{Salida:} Transformación homogénea ${}^{0}T_6$[cite: 1].
    \end{itemize}
    
    \begin{alertblock}{Atención Conceptual}
        Si tengo ${}^{0}T_6$, calcular los $q_i$ corresponde a la cinemática \textbf{inversa}, no directa[cite: 1, 2].
    \end{alertblock}
\end{frame}

\begin{frame}{De una Fila D-H a una Transformación Homogénea}
    \textbf{Convención D-H Estándar de la Clase:}
    \begin{equation*}
        \boxed{ {}^{i-1}T_i = R_z(\theta_i) T_z(d_i) T_x(a_i) R_x(\alpha_i) } \text{[cite: 1]}
    \end{equation*}

    El desarrollo matricial arroja[cite: 1]:
    \begin{equation*}
        {}^{i-1}T_i =
        \begin{bmatrix}
        c_{\theta_i} & -s_{\theta_i}c_{\alpha_i} & s_{\theta_i}s_{\alpha_i} & a_i c_{\theta_i} \\
        s_{\theta_i} & c_{\theta_i}c_{\alpha_i} & -c_{\theta_i}s_{\alpha_i} & a_i s_{\theta_i} \\
        0 & s_{\alpha_i} & c_{\alpha_i} & d_i \\
        0 & 0 & 0 & 1
        \end{bmatrix}
    \end{equation*}

    \textbf{Importante (Offsets articulares):} Para una junta revoluta, $\theta_i = q_i + \theta_{i,0}$[cite: 1]. El cero del fabricante y el cero D-H no siempre coinciden[cite: 1].
\end{frame}

\begin{frame}{Composición de Transformaciones}
    Para determinar la pose global, las transformaciones consecutivas se componen mediante multiplicación matricial en un orden específico[cite: 1, 2].

    \begin{center}
    \begin{tikzpicture}[
        node distance=1.5cm,
        frame/.style={circle, draw=ucbblue, thick, inner sep=2pt, font=\small},
        arrow/.style={-Stealth, thick, ucbblue}
    ]
        \node[frame] (f0) {$\{0\}$};
        \node[frame, right=of f0] (f1) {$\{1\}$};
        \node[frame, right=of f1] (f2) {$\{2\}$};
        \node[right=of f2] (dots) {$\cdots$};
        \node[frame, right=of dots] (f6) {$\{6\}$};

        \draw[arrow] (f0) -- node[above]{\footnotesize ${}^{0}T_1$} (f1);
        \draw[arrow] (f1) -- node[above]{\footnotesize ${}^{1}T_2$} (f2);
        \draw[arrow] (f2) -- node[above]{\footnotesize ${}^{2}T_3$} (dots);
        \draw[arrow] (dots) -- node[above]{\footnotesize ${}^{5}T_6$} (f6);
    \end{tikzpicture}
    \end{center}

    \vspace{0.2cm}
    \textbf{Ecuación de la cadena:}
    \begin{equation*}
        \boxed{ {}^{0}T_6 = {}^{0}T_1 {}^{1}T_2 {}^{2}T_3 {}^{3}T_4 {}^{4}T_5 {}^{5}T_6 } \text{[cite: 1]}
    \end{equation*}
    \textit{Error común: Revertir el orden de multiplicación arruina la semántica de los marcos[cite: 1].}
\end{frame}

\begin{frame}{Interpretación de la Pose Resultante}
    El resultado de la cinemática directa codifica posición y orientación[cite: 1]:
    
    \begin{equation*}
        {}^{0}T_6 =
        \begin{bmatrix}
        \mathbf{n} & \mathbf{o} & \mathbf{a} & \mathbf{p} \\
        0 & 0 & 0 & 1
        \end{bmatrix}
        =
        \begin{bmatrix}
        {}^{0}R_6 & {}^{0}\mathbf{p}_6 \\
        \mathbf{0} & 1
        \end{bmatrix}
    \end{equation*}

    \begin{itemize}
        \item ${}^{0}\mathbf{p}_6$: Vector de posición (origen del marco final en la base)[cite: 1].
        \item ${}^{0}R_6$: Matriz de orientación (vectores direccionales $x_6, y_6, z_6$ expresados en la base)[cite: 1].
    \end{itemize}

    \textbf{Verificación Estructural[cite: 1]:}
    Revisar siempre que $R^T R = I$ y $\det(R) = 1$.
\end{frame}

\begin{frame}{Ejemplo Guiado: Manipulador 2R Planar}
    \begin{columns}
        \begin{column}{0.4\textwidth}
            \textbf{Datos[cite: 1]:} \\
            $a_1 = 0.40\text{ m}, \quad a_2 = 0.30\text{ m}$ \\
            $q_1 = 30^\circ, \quad q_2 = 60^\circ$ \\
            
            \vspace{0.4cm}
            \begin{tikzpicture}[scale=1.5, thick]
                \draw[->, gray] (0,0) -- (1,0) node[right]{$x_0$};
                \draw[->, gray] (0,0) -- (0,1) node[above]{$y_0$};
                
                \coordinate (O) at (0,0);
                \coordinate (J1) at (30:0.4);
                \coordinate (TCP) at ($(J1) + (90:0.3)$);
                
                \draw[ucbblue, line width=2pt] (O) -- (J1) node[midway, below right]{\footnotesize $a_1$};
                \draw[ucbblue!70, line width=2pt] (J1) -- (TCP) node[midway, right]{\footnotesize $a_2$};
                \filldraw[black] (O) circle (1.5pt);
                \filldraw[black] (J1) circle (1.5pt);
                \filldraw[red] (TCP) circle (1.5pt) node[above]{\footnotesize TCP};
            \end{tikzpicture}
        \end{column}
        
        \begin{column}{0.6\textwidth}
            \textbf{Matrices Individuales:}
            \footnotesize
            \begin{equation*}
                {}^{0}T_1 = \begin{bmatrix}
                c_1 & -s_1 & 0 & a_1 c_1 \\
                s_1 & c_1 & 0 & a_1 s_1 \\
                0 & 0 & 1 & 0 \\
                0 & 0 & 0 & 1
                \end{bmatrix}, \quad
                {}^{1}T_2 = \begin{bmatrix}
                c_2 & -s_2 & 0 & a_2 c_2 \\
                s_2 & c_2 & 0 & a_2 s_2 \\
                0 & 0 & 1 & 0 \\
                0 & 0 & 0 & 1
                \end{bmatrix}
            \end{equation*}
            \normalsize
            
            \textbf{Cadena:}
            \footnotesize
            \begin{equation*}
                {}^{0}T_2 = {}^{0}T_1 {}^{1}T_2 = \begin{bmatrix}
                c_{12} & -s_{12} & 0 & a_1 c_1 + a_2 c_{12} \\
                s_{12} & c_{12} & 0 & a_1 s_1 + a_2 s_{12} \\
                0 & 0 & 1 & 0 \\
                0 & 0 & 0 & 1
                \end{bmatrix}
            \end{equation*}
        \end{column}
    \end{columns}
\end{frame}

\begin{frame}{Ejemplo Guiado: Resultado 2R Numérico}
    Sustituyendo $q_1 = 30^\circ, q_2 = 60^\circ \implies q_1+q_2 = 90^\circ$[cite: 1]:
    
    \begin{equation*}
        \boxed{
        {}^{0}T_2 \approx
        \begin{bmatrix}
        0 & -1 & 0 & 0.3464 \\
        1 & 0 & 0 & 0.5000 \\
        0 & 0 & 1 & 0 \\
        0 & 0 & 0 & 1
        \end{bmatrix}
        } \text{[cite: 1]}
    \end{equation*}

    \begin{itemize}
        \item \textbf{Posición:} ${}^{0}\mathbf{p}_2 = [0.3464, 0.5000, 0]^T$ m[cite: 1].
        \item \textbf{Orientación:} Eje $x_2$ apunta hacia $y_0$ (Rotación global de $90^\circ$).
        \item \textbf{Verificación geométrica:} $\sqrt{0.3464^2 + 0.5000^2} \approx 0.608 < 0.70$ m (Alcance máximo $a_1+a_2$). ¡Físicamente posible![cite: 1].
    \end{itemize}
\end{frame}

\begin{frame}{Aplicación ABB IRB 120: Tabla D-H Publicada}
    \textbf{Parametrización reportada por Truc \& Lam (2020)[cite: 1]:}
    
    \begin{table}[]
        \centering
        \renewcommand{\arraystretch}{1.2}
        \begin{tabular}{c c c c c}
        \toprule
        $i$ & $\theta_i$ & $d_i$ [m] & $a_i$ [m] & $\alpha_i$ \\ \midrule
        1 & $q_1$ & 0.290 & 0 & $-\pi/2$ \\
        2 & $q_2 - \pi/2$ & 0 & 0.270 & 0 \\
        3 & $q_3$ & 0 & 0.070 & $-\pi/2$ \\
        4 & $q_4$ & 0.302 & 0 & $+\pi/2$ \\
        5 & $q_5$ & 0 & 0 & $-\pi/2$ \\
        6 & $q_6$ & 0.072 & 0 & 0 \\ \bottomrule
        \end{tabular}
    \end{table}
    
    \begin{alertblock}{Cero Articular}
        Observen que si $q_2 = 0$, el ángulo $\theta_2$ efectivo para la transformación no es cero, sino $-\pi/2$[cite: 1].
    \end{alertblock}
\end{frame}

\begin{frame}{Aplicación ABB IRB 120: Evaluación en Configuración Cero}
    Para la configuración didáctica $\mathbf{q} = [0, 0, 0, 0, 0, 0]^T$[cite: 1]:

    \begin{columns}
        \begin{column}{0.5\textwidth}
            \footnotesize
            $ {}^{0}T_1 = \begin{bmatrix} 1&0&0&0 \\ 0&0&1&0 \\ 0&-1&0&0.290 \\ 0&0&0&1 \end{bmatrix}$ \\ \vspace{0.1cm}
            $ {}^{1}T_2 = \begin{bmatrix} 0&1&0&0 \\ -1&0&0&-0.270 \\ 0&0&1&0 \\ 0&0&0&1 \end{bmatrix}$ \\ \vspace{0.1cm}
            $ {}^{2}T_3 = \begin{bmatrix} 1&0&0&0.070 \\ 0&0&1&0 \\ 0&-1&0&0 \\ 0&0&0&1 \end{bmatrix}$
        \end{column}
        
        \begin{column}{0.5\textwidth}
            \footnotesize
            $ {}^{3}T_4 = \begin{bmatrix} 1&0&0&0 \\ 0&0&-1&0 \\ 0&1&0&0.302 \\ 0&0&0&1 \end{bmatrix}$ \\ \vspace{0.1cm}
            $ {}^{4}T_5 = \begin{bmatrix} 1&0&0&0 \\ 0&0&1&0 \\ 0&-1&0&0 \\ 0&0&0&1 \end{bmatrix}$ \\ \vspace{0.1cm}
            $ {}^{5}T_6 = \begin{bmatrix} 1&0&0&0 \\ 0&1&0&0 \\ 0&0&1&0.072 \\ 0&0&0&1 \end{bmatrix}$
        \end{column}
    \end{columns}
\end{frame}

\begin{frame}{Aplicación ABB IRB 120: Extracción de Pose}
    Multiplicando la cadena ${}^{0}T_6 = {}^{0}T_1 {}^{1}T_2 \cdots {}^{5}T_6$ numéricamente para $\mathbf{q}=\mathbf{0}$[cite: 1]:
    
    \begin{equation*}
        \boxed{
        {}^{0}T_6 =
        \begin{bmatrix}
        0 & 0 & 1 & 0.374 \\
        0 & -1 & 0 & 0 \\
        1 & 0 & 0 & 0.630 \\
        0 & 0 & 0 & 1
        \end{bmatrix}
        } \text{[cite: 1]}
    \end{equation*}

    \begin{itemize}
        \item \textbf{Posición Base-a-Brida:} ${}^{0}\mathbf{p}_6 = [0.374, \; 0, \; 0.630]^T \text{ m}$[cite: 1].
        \item \textbf{Ejes de la Brida:} $x_6 \parallel +z_0$, $y_6 \parallel -y_0$, $z_6 \parallel +x_0$[cite: 1].
    \end{itemize}
\end{frame}

\begin{frame}{Advertencia Técnica: Convenciones de Offset}
    \begin{tcolorbox}[colback=white, colframe=ucbblue, title=Modelos Públicos Discrepantes para el IRB 120]
        Dos fuentes científicas modelan la misma geometría del IRB 120, pero difieren en el "cero" de la junta 6[cite: 1]:
        
        \begin{itemize}
            \item \textbf{Modelo A (Truc \& Lam):} $\theta_6 = q_6$[cite: 1]
            \item \textbf{Modelo B (Bahani et al.):} $\theta_6 = q_6 + \pi$[cite: 1]
        \end{itemize}
        
        \textbf{Lección crítica:} Mismo robot físico $\neq$ misma parametrización algebraica[cite: 1]. Si mezclan valores de fuentes distintas sin homogeneizar las referencias de marco, la orientación en la simulación fallará[cite: 1].
    \end{tcolorbox}
\end{frame}

\begin{frame}{Comprobación Formativa (Exit Ticket)}
    \begin{tcolorbox}[colback=white, colframe=ucbblue, title=Consolidación Rápida]
        \begin{enumerate}
            \item En cinemática directa, las variables articulares son \rule{2cm}{0.15mm} y la pose del efector final es \rule{2cm}{0.15mm}[cite: 1, 2].
            \item Si se conocen ${}^{0}T_1$, ${}^{1}T_2$ y ${}^{2}T_3$, escriba la ecuación para ${}^{0}T_3$[cite: 1, 2].
            \item Dada la siguiente matriz, identifique el vector de posición[cite: 1, 2]:
            $T = \begin{bmatrix} 0&-1&0&0.20 \\ 1&0&0&0.35 \\ 0&0&1&0 \\ 0&0&0&1 \end{bmatrix}$
            \item Si la tabla D-H estipula $\theta_2 = q_2 - \pi/2$, ¿cuál es el ángulo $\theta_2$ que debe ingresar a la matriz si el usuario comanda $q_2 = 0$?[cite: 1]
        \end{enumerate}
    \end{tcolorbox}
\end{frame}

\end{document}
